% ═══════════════════════════════════════════════════════════════════
% ЧАСТЬ XIV. ТЕОРЕТИКО-ГРУППОВЫЕ СЛОИ
% ═══════════════════════════════════════════════════════════════════
\part{Теоретико-групповые слои программы}

\section{Группа $\Gamma(2,3,7)$ и её роли}\label{sec:gamma237}

Треугольная группа $\Gamma(2,3,7)$ --- группа автоморфизмов
квартики Клейна порядка $336$ --- играет в программе три
комплементарные роли, и полезно различить их явно.

\emph{Роль геометрическая.} $\Gamma(2,3,7)$ действует на кривой
Клейна перестановками, сохраняющими все конструкции; её элементы
порядков $2,3,7$ поставляют фазы $\pi/2$, $\pi/3$, $\pi/7$
(семейство $\pi/N$). Образующая $C$ порядка $7$ --- источник
главной фазы $\delta_C=\pi/7$ всей программы.

\emph{Роль арифметическая.} Группа реализуется как
$\mathrm{PSL}(2,7)$ над конечным полем; её представительства
связывают гептаду $\sqrt{-7}$ с кодом Фано и с дискриминантом
$\Delta=-7^3$. Арифметика порядков --- точная: $168=2^3\cdot3\cdot7$,
и оба разложения ($168$ как $\mathrm{PSL}(2,7)$ и как группа
треугольника) читаются в одной таблице характеров.

\emph{Роль конвейерная.} Через $\mu_4\subset\Gamma(2,3,7)$
программа связывает фазу $\pi/7$ с эквивариантностью тора:
сертификаты~C и~D показывают, что одна и та же $\mu_4$-механика
работает на всех стендов. Это позволяет переносить сертификаты
между ступенями: доказанная на торе коммутация $[L,P]=0$
переиздаётся на K3 и на Клейне без изменений доказательства.

\begin{proposition}[три порядка --- три фазы]\label{prop:threeorders}
Элементы $\Gamma(2,3,7)$ порядков $2,3,7$ дают в спин-двойном
накрытии фазы $\pi/2$, $\pi/3$, $\pi/7$ соответственно; все три
фазы реализуются в периодах якобиана: характеры $\mu_2$, $\mu_3$,
$\mu_7$ присутствуют в разложении $H^0(\Omega^1)$ и в
период-отношениях.
\end{proposition}

\begin{proof}
Фаза $\pi/N$ --- общее правило \eqref{eq:delta} для элемента
порядка $N$. Реализация: изотипика $(1,1,0,1)$ по $\mu_4$
(теорема~\ref{th:klein}) показывает механизм характеров на
$H^0(\Omega^1)$; тождество характера $\zeta+\zeta^2+\zeta^4=
\tau-1$ реализует $\mu_7$ в период-отношении; $\mu_2$ и $\mu_3$
реализуются соответствующими подгруппами в той же таблице
характеров. Численные период-отношения согласованы с точными
значениями до $10^{-30}$ (протокол Клейна, проверки K2).
\end{proof}

\section{Перестановочные представления и их спектры}\label{sec:permspec}

Сертификаты C--D используют одну технику --- спектр
перестановочного представления. Изложим её в общем виде.

Пусть конечная группа $G$ действует на конечном множестве
$\Omega$ из $m$ точек. Перестановочное представление
$\rho\colon G\to\mathrm{GL}(\CC^m)$ раскладывается по неприводимым
характерам; кратности вычисляются проекционными операторами
\[
  \Pi_\chi \;=\; \frac{\dim\chi}{|G|}
  \sum_{g\in G}\overline{\chi(g)}\,\rho(g).
\]
Кратность характера $\chi$ равна $\rank\Pi_\chi$ --- целому числу,
вычисляемому точной арифметикой, если матрицы $\rho(g)$
целочисленные. В программе:
\begin{itemize}[leftmargin=1.6em, itemsep=0.1em]
\item $G=\mu_4$, $\Omega=48$ прямых K3: кратности $(1,7,7,7)$
  (теорема~\ref{th:C1}); три пути подсчёта согласованы;
\item $G=\mu_4$, $\Omega=$ решётка тора: регулярное представление,
  $x^4-1$ (теорема~\ref{th:C2});
\item $G=\Gamma(2,3,7)$, $\Omega=7$ точек Фано: стандартное
  $3$-мерное представление и гептада (пропозиция~\ref{prop:fano});
\item $G=(\mu_4)^2$, $\Omega=48$: двенадцать орбит, кратности
  согласованы с инерцией.
\end{itemize}
Везде спектр --- целочисленный объект, и его согласование с
независимыми путями (прямые орбиты, проекторы, характеры) ---
сертификат первого типа.

\section{Сравнение ступеней: параллельная таблица}\label{sec:compare}

Сведём ступени в параллельную таблицу по всем слоям конвейера ---
это позволяет увидеть единообразие правил при росте масштаба.

\begin{center}
\small
\begin{tabularx}{\textwidth}{@{}>{\raggedright\arraybackslash}p{2.5cm}
  >{\raggedright\arraybackslash}X
  >{\raggedright\arraybackslash}X
  >{\raggedright\arraybackslash}X@{}}
\toprule
Слой & Тор & Клейн & Ферма 15/30 \\
\midrule
Тройка & $(4,1,4\pi^2)$ & $(7,22,3.338)$ &
  $(15,\cdot,\cdot)$, $(30,\cdot,\cdot)$ \\
Фазы & $\pi/4$ & $\pi/2,\pi/3,\pi/7$ & $\pi/15$, $\pi/30$ \\
Поля решёток & $\ZZ[i]$ & $\QQ(\sqrt{-7})$ &
  $\QQ(\zeta_{15})=\QQ(\zeta_{30})$ \\
Дискриминанты & $1$ & $7^3$ & $1125^2$ \\
Эквивариантность & $\mu_4$ точная & $\mu_4$+гептада &
  $\mu_N\times\mu_N$ точная \\
SNF-тип & $(1)$ & $(1,1,1)$ &
  $(1,1,5,5,15,15,15,15)$ \\
Сертификаты & вывод целиком & B, C & B(H3), C(H3), F, G, H \\
\bottomrule
\end{tabularx}
\end{center}

Таблица показывает главный принцип: ни одна строка не является
«специальной настройкой» под ступень --- каждая строка заполнена
одним и тем же правилом, применённым к разным тройкам. Рост
дискриминантов ($1\to7^3\to1125^2$) отражает рост арифметической
богатства полей, и конвейер отслеживает его без изменения
алгоритмов --- SNF-модуль, перепись, LLL-слой работают на всех
уровнях одинаково.

\subsection{Роль K3 как универсального носителя}

K3-слой присутствует на всех ступенях через индекс $b_2=22$: на
торе он даёт альтернативный носитель $k=22$ (таблица
\ref{tab:torus-triples}), на Клейне --- основной носитель тройки,
на верхних ступенях --- наследуется через сертификат~G5
(квадрант-интеграл $\Ga(1/4)^4/(16\pi\sqrt2)$). Универсальность
K3 --- конструктивный факт: решётка K3 унимодулярна сигнатуры
$(3,19)$, и её второе число Бетти $22$ --- наименьший
целочисленный индекс, общий для всех ступеней лестницы. Именно
поэтому фреймворк фиксирует $k=22$ как «универсальный выбор» ---
это часть определения конвейера, а не подбор.

\section{Развёрнутый разбор отчёта лаборатории}\label{sec:reportwalk}

Запуск «Полный протокол» создаёт в \texttt{reports/} комплект
файлов. Разберём структуру отчёта --- это карта того, как
лаборатория кодирует результаты монографии.

\subsection{Схема JSON-отчёта}

\begin{verbatim}
{
  "meta": {"version": ..., "dps": 35, "language": "ru|en"},
  "V1_census":  {"15": {"h_d": {...}, "genus": 91, "total": 91},
                 "30": {...}},
  "V2_certB":   {"15": {"n_tests": 21, "max_rel_err": ...,
                        "per_conductor": {...}}, "30": {...}},
  "V3_phase":   {"15": {"max_phase_dev_rad": 0.0}, ...},
  "V4_certC":   {"max_rel_err": ...},
  "V5_chain":   {"max_abs": ...},
  "V6_reflection": {"max_rel_err": ...},
  "V7_rank":    {"15": {"rank": 91, "cond": 8.63}, ...},
  "V8_orthogonal": {"diag": ..., "offdiag": ...},
  "V9_deep":    {"rel_err": ...},
  "stands":     {"torus": {...}, "k3": {...}, "klein": {...},
                 "errata": {...}, "binary_code": {...}},
  "certificates": {"A": {...}, ..., "H": {...}},
  "verdict":    {"all_pass": true, "thresholds": {...}}
}
\end{verbatim}

Каждый ключ соответствует разделу монографии
(соответствие --- таблица~\ref{tab:vprotocol} и
раздел~\ref{sec:labmap}); вердикт \texttt{all\_pass}
вычисляется по порогам автоматически. Схема стабильна:
эталонные файлы \texttt{results/} сверяются с новыми отчётами
по ключам и порогам, а не по битовым совпадениям --- это
позволяет воспроизводить на других платформах.

\subsection{Текстовый журнал}

Журнал \texttt{verification\_log.txt} --- человекочитаемая версия
отчёта: каждая проверка печатает заголовок, действия, числа и
вердикт PASS. Журнал предназначен для быстрого просмотра и для
вставки в задачи воспроизведения; JSON --- для автоматических
сверок.

\subsection{Плиточные диаграммы}

Каталог \texttt{reports/plots/} содержит восемь плиток $600$~dpi:
шесть плиток атласа (раздел~\ref{sec:atlas}) плюс плитки
конкретного запуска (спектр ранга V7, распределения остатков V2 по
проводникам). Плиточная система --- модульная сетка: каждая плитка
строится независимо, подписывается метаданными запуска
($dps$, дата, параметры) и не зависит от остальных --- это
позволяет добавлять новые плитки без перекомпиляции старых.
